\begin{prob}
    % Write the problem here.
    Suppose a professor makes the following deal with their class: by attending
    lecture, you can earn ``tokens''. At the end of the quarter, each token can either
    be used to flip a late homework to on-time, or it can be redeemed for extra credit.

    Suppose the first token that is redeemed is awarded 50 points of extra
    credit, the second 30 points, the third 20 points, and the fourth 10
    points. Beyond this, all tokens are worth 10 points -- that is, the fifth,
    sixth, seventh, etc., tokens are all worth 10 points if redeemed. For
    example, if a student has three tokens and uses the first to flip a late
    homework and redeems the remaining two for extra credit, they'll earn 50 +
    30 = 80 points of extra credit. Also suppose that any late homework (which
    isn't flipped) is given zero points.

    If the goal is to maximize the number of points earned, it might seem like
    the most advantageous strategy for the student is to first use the tokens
    to flip all of their late homeworks and only then redeem their remaining
    tokens for extra credit, because late homeworks would receive zero points
    otherwise. But suppose a student's late homework earned only 1 point --
    clearly, they shouldn't waste a token on this homework, and should instead
    redeem it for extra credit. That is, the optimal strategy requires a little
    more thought.

    In a file named \python{allocate_tokens.py}, write a function
    \python{allocate_tokens(lates, n_tokens)} which takes in a list
    \python{lates} of the points earned by a student on all of their late
    homeworks (in arbitrary order), as well as \python{n_tokens}, the number of tokens they have
    earned. It should return the total number of points earned by using the
    tokens in the best possible way (that is, the maximum number of points
    possible).

    For example, suppose a student's two late homeworks earned 75 and 20 points, respectively, and
    they earned two tokens.
    Then \python{lates = [20, 75]}, and \python{n_tokens = 2}. The best strategy is for them to use
    one token to ``flip'' the 75 point homework to on-time, and to redeem the second token
    for 50 points of extra credit, earning a total of 125 points. Therefore, \python{allocate_tokens([20, 75], 2)}
    should return \python{125}.

    Your code will still receive partial credit if it is inefficient, but to receive full
    credit it should have a worst-case time complexity of $\Theta(n \log n)$,
    where $n$ is the number of late homeworks.\footnote{The time taken will also depend on
    the number of tokens earned, but you can assume that the number of tokens earned is $O(n)$.}

    \emph{Hint}: the worst-case time complexity of $\Theta(n \log n)$ should point you towards
    a strategy. What else do we know of that takes $\Theta(n \log n)$ time? Also be sure that your code works when
    there are fewer late assignments than tokens earned. It is a good idea to write several of your
    own doctests to thoroughly check your code.


    \begin{soln}
        The fact that our worst-case time complexity should be $\Theta(n \log n)$ suggests that maybe
        \emph{sorting} the input would be useful. What if we sort \python{lates}; does that help?

        It turns out that it does -- \emph{if} we sort from largest to smallest. The strategy
        is as follows: After sorting \python{lates} in descending order, we loop through and
        compare the number of points earned on that assignment to the number of points we'd
        earn by redeeming a token (if we have any remaining). If the token is not worth
        more in extra credit than the late assignment, we ``flip'' the late assignment and move
        on to the next late. However, if the token \emph{is} worth more in extra credit, we redeem the token for
        extra credit. In this case, we do not move on to the next late assignment -- we might want
        to use the next token to flip this late.

        Note that this looping approach does not work if we don't sort the lates first. If we
        looped through the late assignments in arbitrary order, there could be a very valuable
        late lurking at the end of the list for which we should use a token. Without sorting,
        we'd need to find the largest remaining late on every iteration, leading to a
        quadratic overall time complexity. As-is, the code takes $\Theta(n \log n)$ because
        of the sorting. All of the code after the sort takes linear time.

        Here is what a solution looks like:

        \inputminted{python}{\thisdir/include-solution/allocate_tokens.py}

    \end{soln}

\end{prob}
